%
% Polynome von Erweiterungen
%
\subsection{Polynome von Erweiterungen
\label{buch:intalg:liouville:subsection:polynome}}
Die oben formulierte Beweisstrategie für das Prinzip von Liouville
basiert darauf, dass wir die Ableitung der Stammfunktion ausrechnen
und dann nachweisen, dass nur logarithmische Terme als Ableitung
Ausdrücke in $F$ ergeben.
Dazu müssen zunächst die Ableitungen von Zähler und Nenner in
\eqref{buch:intalg:liouville:eqn:bruch}
verstanden werden, was in den folgenden Abschnitten geklärt werden
soll.

%
% Polynome einer logarithmischen Erweiterung
%
\subsubsection{Polynome einer logarithmischen Erweiterung}
Das Prinzip von Liouville verspricht, dass nur logarithmische Erweiterungen
für die Stammfunktion gebraucht werden, und dass diese nur linear
vorkommen.
Dies wird sich darin äussern, dass die Erweiterung im rationalen Teil
nicht vorkommen kann.
Dies wird sich daraus ergeben, wie genau sich der Grad eines Polynoms
in $\vartheta$ beim Ableiten ändert.
Der folgende Satz gibt darüber Auskunft.

\begin{satz}
\label{buch:intalg:liouville:polynome:satz:logarithmisch}
Sei $F$ ein Funktionenkörper mit Konstantenkörper $K\subset F$ und sei
$\vartheta$ ein logarithmisch transzendentes Element über $F$ mit
$\vartheta=\log u$ mit $u\in F$.
Sei ausserdem
$a(\vartheta) = a_0+a_1\vartheta+\dots+a_n\vartheta^n\in F[\vartheta]$
ein Polynom in $\vartheta$ mit Koeffizienten in $F$.
Dann gilt:
\begin{enumerate}
\item
Die Ableitung von $a(\vartheta)$ ist ebenfalls ein Polynom mit Koeffizienten
in $F$: $Da(\vartheta)\in F[\vartheta]$.
\item
Wenn der Leitkoeffizient $a_n$ eine Konstante ist, dann ist
$\deg Da(\vartheta)=\deg a(\vartheta) -1$.
\item
Wenn der Leitkoeffizient $a_n$ keine Konstante ist, dann ist
$\deg Da(\vartheta)=\deg a(\vartheta)$.
\end{enumerate}
\end{satz}

\begin{proof}
\begin{enumerate}
\item
Die Ableitung ist
\begin{align*}
Da(\vartheta)
&=
D\sum_{k=0}^n a_k\vartheta^k
=
\sum_{k=0}^n \bigl(Da_k\vartheta^k + a_kk\vartheta^{k-1}D\vartheta\bigr)
\\
&=
\sum_{k=0}^n Da_k\vartheta^k
+
D\vartheta \sum_{k=0}^{n-1} a_{k+1}(k+1)\vartheta^k
\\
&=
\sum_{k=0}^{n-1} \bigl(Da_k+(k+1)a_{k+1}D\vartheta\bigr) \vartheta^k
+
Da_n \vartheta^n
\\
&=
\sum_{k=0}^{n-1}
\biggl(
\underbrace{Da_k+(k+1)a_{k+1}\frac{Du}{u}}_{\displaystyle\in F}
\biggr) \vartheta^k
+
Da_n \vartheta^n
\end{align*}
Da $Du\in F$ und $u\in F$ folgt, dass der Koeffizient in der Klammer
ebenfalls in $F$ ist und daher ist $Da(\vartheta)\in F[\vartheta]$.
\item
Wenn $a_n$ eine Konstante ist, dann fällt der Term mit $\vartheta^n$ weg,
der Grad von $Da(\vartheta)$ ist also höchstens $n-1$.
Es ist zu zeigen, dass der Grad genau $n-1$ ist, d.~h.~dass der Koeffizient
von $\vartheta^{n-1}$ nicht verschwindet.
Würde dieser Koeffizient verschwinden, dann wäre wegen $Da_n=0$
\begin{align*}
0
&=
Da_{n-1} + na_n D\vartheta
\\
&=
D(a_{n-1}+na_n\vartheta),
\end{align*}
d.~h.~$a_{n-1}+na_n\vartheta=c$ ist eine Konstante $c\in K$.
Die Gleichung
\[
na_n\cdot\vartheta
+
(\underbrace{a_{n-1}-c}_{\displaystyle\in K})
=
0
\]
bedeutet aber, dass $\vartheta$ algebraisch ist über $K$.
Dies widerspricht der Voraussetzung, dass $\vartheta$ transzendent ist.
Der Widerspruch zeigt, dass $\deg Da(\vartheta)=n-1$ ist.
\item
Wenn $a_n$ keine Konstante ist, dann ist $Da_n\ne 0$ und damit
$\deg Da(\vartheta)=n$.
\qedhere
\end{enumerate}
\end{proof}

%
% Polynome einer exponentiellen Erweiterung
%
\subsubsection{Polynome einer exponentiellen Erweiterung}
Im Gegensatz zu einer logarithmischen Erweiterung nimmt der Grad
eines Polynoms einer exponentiellen Erweiterung beim Ableiten nicht ab.

\begin{satz}
\label{buch:intalg:liouville:polynome:satz:exponentiell}
Sei $F$ ein Funktionenkörper mit Konstantenkörper $K\subset F$ und sei
$\vartheta$ ein exponentielles transzendentes Element über $F$ mit
$\vartheta=\exp u$ mit $u\in F$.
Sei ausserdem
$a(\vartheta) = a_0+a_1\vartheta+\dots+a_n\vartheta^n\in F[\vartheta]$
ein Polynom in $\vartheta$ mit Koeffizienten in $F$.
Dann gilt:
\begin{enumerate}
\item
Für $h\in F$ und $n\ne 0$ gilt $D(h\vartheta^n) = \bar{h}\vartheta^n$ mit
$\bar{h}\in F$ und $\bar{h}\ne 0$.
\item
Für ein Polynom
$a(\vartheta)\in F[\vartheta]$
gilt
$\deg Da(\vartheta) = \deg a(\vartheta)$.
\item
$a(\vartheta) \in F[\vartheta]$
teilt $Da(\vartheta)$ genau dann, wenn $a(\vartheta)$ ein Monom ist,
also $a(\vartheta)=h\vartheta^n$ mit $h\in F\setminus\{0\}$, $n\in\mathbb{Z}$.
\end{enumerate}
\end{satz}

\begin{proof}
\begin{enumerate}
\item
Die Ableitung ist
\[
D(h\vartheta^n)
=
Dh\cdot\vartheta^n + hn\vartheta^{n-1}D\vartheta
=
Dh\cdot\vartheta^n + hn\vartheta^{n-1} \vartheta Du
=
(\underbrace{Dh+ hn Du}_{\displaystyle=\bar{h}}) \vartheta^n
\]
Wäre $\bar{h}=0$, dann wäre $h\vartheta^n$ eine Konstante $c\in K$,
d.~h.~es gälte eine Gleichung der Form
\[
h\cdot \vartheta^n - c = 0.
\]
Dies bedeutet, dass $\vartheta$ algebraisch ist über $F$, im Widerspruch
zur Voraussetzung, dass $\vartheta$ transzendent ist.
Somit folgt $\bar{h}\ne 0$.
\item
Die Ableitung des Terms $a_n\vartheta^n$ des Polynoms ist
$\bar{a}_n\vartheta^n$ mit $\bar{a}_n\ne 0$.
Somit ist $\deg Da(\vartheta) = n = \deg a(\vartheta)$.
\item
Wenn $a(\vartheta)$ das Monom $h\vartheta^n$ ist, dann ist
\[
Da(\vartheta)
=
\bar{h}\vartheta^n
=
\frac{\bar{h}}{h}
\cdot
h\vartheta^n
=
\frac{\bar{h}}{h}
\cdot
a(\vartheta),
\]
d.~h.~wegen $\bar{h}/h\in F$ ist $Da(\vartheta)$ ist ein Teiler von
$a(\vartheta)$.

Sei jetzt umgekehrt $Da(\vartheta)$ ein Teiler von $a(\vartheta)$.
Es gibt dann ein $g\in F$ mit $Da(\vartheta)=ga(\vartheta)$.
Sei $m$ der höchste Exponent $<n$ von $\vartheta$ im Polynom $a(\vartheta)$,
welcher einen nicht verschwindenen Koeffizienten $a_m\ne 0$ hat.
Das Polynom hat also die Form
\begin{align*}
a(\vartheta)
&=
a_n\vartheta^n + a_m\vartheta^m + \ldots + a_1\vartheta + a_0
\intertext{mit der Ableitung}
Da(\vartheta)
&=
\bar{a}_n\vartheta^n
+
\bar{a}_m\vartheta^m
+
\ldots
+
\bar{a}_1\vartheta
+
Da_0.
\end{align*}
mit $\bar{a}_k=Da_k+ka_kDu$.
Die Teilbarkeitsbedingung bedeutet dann
\[
\renewcommand{\arraycolsep}{2pt}
Da(\vartheta)
=
ga(\vartheta)
\quad
\Rightarrow
\quad
\left\{\;
\begin{array}{lclclcl}
ga_n &=& \bar{a}_n &=& Da_n &+& na_n Du \\
ga_m &=& \bar{a}_m &=& Da_m &+& ma_m Du.
\end{array}
\right.
\]
Auflösen nach $g$ ergibt
\begin{align*}
g
&=
\frac{Da_n}{a_n} + n\,Du
=
\frac{Da_m}{a_m} + m\,Du.
\intertext{Die Differenz der rechten Seiten ist}
0
&=
\frac{Da_n}{a_n}
-
\frac{Da_m}{a_m}
+
(n-m)\,Du.
%\intertext{Die rechte Seite kann man auch als}
%0
%&=
%D
%\biggl(
%\frac{a_n}{a_m}
%+
%(n-m)u
%\biggr)
\end{align*}
Die rechte Seite ermöglicht auch, die Ableitung
\begin{align*}
D
\biggl(
\frac{a_n}{a_m}\vartheta^{n-m}
\biggr)
&=
\frac{(Da_n) a_m - a_n\,Da_m}{a_m^2} \vartheta^{n-m}
+
(n-m)\vartheta^{n-m}Du
\\
&=
\frac{a_n}{a_m}
\biggl(
\frac{Da_n}{a_n}
-
\frac{Da_m}{a_m}
-
(n-m)Du
\biggr)
=
0
\end{align*}
zu berechnen.
Somit ist $(a_n/a_m)\vartheta^{n-m}$ eine Konstante, was nur für $n-m=0$
möglich ist.
Es gibt also gar keinen Koeffizienten $a_m\ne 0$ mit $m<n$.
Somit ist $a(\vartheta)=a_n\vartheta^n$ ein Monom.
\qedhere
\end{enumerate}
\end{proof}

%
% Polynome einer algebraischen Erweiterung
%
\subsubsection{Polynome einer algebraischen Erweiterung}
Logarithmische und exponentielle transzendente Erweiterungen
sind durch ihre Eigenschaften unter Differentiation charakterisiert.
Damit ist bereits sichergestellt, dass die Ableitung im erweiterten
Funktionenkörper liegt.
Für eine algebraische Erweiterung $F(\vartheta)$ von $F$ ist dies
jedoch nicht selbstverständlich.
Der folgende Satz zeigt, dass auch für eine algebraische Erweiterung
um das algebraische Element $\vartheta$ die Ableitung $D\vartheta$
in $F(\vartheta)$ liegt.

\begin{satz}
\label{buch:intalg:liouville:polynome:satz:algebraisch}
Sei $\vartheta$ ein algebraisches Element über $F$ mit Minimalpolynom
\begin{equation}
0
=
m(\vartheta)
=
m_0 + m_1\vartheta + \ldots + m_{n-1}\vartheta^{n-1} + \vartheta^n
=
\sum_{k=0}^n m_k\vartheta^k
\label{buch:intalg:liouville:polynome:eqn:minimal}
\end{equation}
mit $m_n=1$.
\begin{enumerate}
\item
Die Ableitung von $\vartheta$ hat die Form
\[
D\vartheta
=
-
\frac{d(\vartheta)}{e(\vartheta}
\in
F(\vartheta)
\]
mit Polynomen 
\begin{equation}
d(\vartheta)
=
\sum_{k=0}^{n-1} (Dm_k) \vartheta^k
\qquad\text{und}\qquad
e(\vartheta)
=
\sum_{k=0}^{n-1} (k+1)m_{k+1} \vartheta^k
=
m'(\vartheta).
\label{buch:intalg:liouville:polynome:eqn:zaehlernenner}
\end{equation}
\item
Eine rationale Funktion von $\vartheta$ mit Koeffizienten in $F$
kann immer als Polynom vom Grad höchstens $n-1$ in $\vartheta$
mit Koeffizienten in $F$ geschrieben werden.
\end{enumerate}
\end{satz}

\begin{proof}
\begin{enumerate}
\item
Ableiten von
\eqref{buch:intalg:liouville:polynome:eqn:minimal}
liefert
\begin{align*}
0
=
Dp(\vartheta)
&=
\sum_{k=0}^n (Dm_k)\vartheta^k
+
\biggl(\sum_{k=0}^n km_k\vartheta^{k-1} \biggl)D\vartheta
\\
&=
\sum_{k=0}^{n-1} (Dm_k)\vartheta^k
+
\biggl(\sum_{k=0}^{n-1} (k+1)m_{k+1}\vartheta^k \biggl)D\vartheta
,
\end{align*}
wobei wir $Dm_n=0$ bereits ausgenützt haben.
Auflösen nach $D\vartheta$ liefert
\[
D\vartheta
=
-\frac{
\displaystyle
\sum_{k=0}^{n-1} (Dm_k) \vartheta^k
}{
\displaystyle
\sum_{k=0}^{n-1} (k+1)m_{k+1} \vartheta^k
}
\in
F(\vartheta).
\]
Zähler und Nenner sind die Polynome in
\eqref{buch:intalg:liouville:polynome:eqn:zaehlernenner}
\item
Sei 
\begin{equation}
f(\vartheta)
=
\frac{u(\vartheta)}{v(\vartheta)}
\label{buch:intalg:liouville:polyalg:eqn:bruch}
\end{equation}
ein Quotient von Polynomen mit Koeffizienten in $F$.
Nach Reduktion mit dem Minimalpolynom darf man annehmen, dass
der Grad von Zähler und Nenner kleiner als $n$ ist.
Insbesondere sind $v(Z)$ und das Minimalpolynom $m(Z)$
teilerfremd.
Nach dem erweiterten euklidischen Algorithmus gibt es
Polynome $s(Z)$ und $t(Z)$ mit 
\[
s(Z)v(Z) + t(Z)m(Z) = 1.
\]
Einsetzen von $\vartheta$ bringt den Term mit $m(\vartheta)=0$ zum
verschwinden und liefert
\[
s(\vartheta)v(\vartheta) = 1
\qquad\Rightarrow\qquad s(\vartheta) = \frac{1}{v(\vartheta)}.
\]
Damit ist das inverse Element von $v(\vartheta)$ als Polynom in
$\vartheta$ ausgedrückt.
Für den Bruch~\eqref{buch:intalg:liouville:polyalg:eqn:bruch}
folgt daher
\[
f(\vartheta)
=
\frac{u(\vartheta)}{v(\vartheta)}
=
u(\vartheta)\cdot s(\vartheta).
\]
Die rechte Seite ist ein Polynom.
Mit dem Minimalpolynom kann es weiter auf ein Polynom vom 
Grad $<n$ reduziert werden.
\qedhere
\end{enumerate}
\end{proof}
